
\documentclass{article}
\usepackage{colortbl}
\usepackage{makecell}
\usepackage{multirow}
\usepackage{supertabular}

\begin{document}

\newcounter{utterance}

\centering \large Interaction Transcript for game `cladder', experiment `full\_v1.5\_default', episode 2017 with score{-}Qwen3.5{-}9B{-}sp.
\vspace{24pt}

{ \footnotesize  \setcounter{utterance}{1}
\setlength{\tabcolsep}{0pt}
\begin{supertabular}{c@{$\;$}|p{.15\linewidth}@{}p{.15\linewidth}p{.15\linewidth}p{.15\linewidth}p{.15\linewidth}p{.15\linewidth}}
    \# & \multicolumn{2}{c}{Player} && \multicolumn{2}{c}{Game Master} \\
    \hline

    \theutterance \stepcounter{utterance}  
    & & & \multicolumn{4}{p{0.6\linewidth}}{
        \cellcolor[rgb]{0.9,0.9,0.9}{
            \makecell[{{p{\linewidth}}}]{
                \texttt{\tiny{[P1$\langle$GM]}}
                \texttt{You are an expert in causal inference. The following question is not a typical commonsense query, but rather a meticulously designed question created by a professor specializing in causal inference, intended to assess the students' mastery of the course content.} \\
\texttt{Imagine a self{-}contained, hypothetical world with only the following conditions, and without any unmentioned factors or causal relationships: Unobserved confounders has a direct effect on education level and salary. Proximity to a college has a direct effect on education level. Education level has a direct effect on salary. Unobserved confounders is unobserved. For people living far from a college, the probability of high salary is 39\%. For people living close to a college, the probability of high salary is 49\%. For people living far from a college, the probability of college degree or higher is 57\%. For people living close to a college, the probability of college degree or higher is 30\%. Will college degree or higher increase the chance of high salary?} \\
\texttt{Start your answer with 'Yes' or 'No', followed by additional reasoning or evidence to support your explanation.} \\
            }
        }
    }
     \\ \\

    \theutterance \stepcounter{utterance}  
    & \multicolumn{4}{p{0.6\linewidth}}{
        \cellcolor[rgb]{0.561,0.737,0.561}{
            \makecell[{{p{\linewidth}}}]{
                \texttt{\tiny{[P1$\rangle$GM]}}
                \texttt{No} \\
\\ 
\texttt{**Reasoning:**} \\
\\ 
\texttt{To determine if a college degree increases the chance of a high salary, we must distinguish between **association** (correlation) and **causation**. In this hypothetical world, the relationship between education and salary is confounded by an unobserved variable, \$U\$ (Unobserved Confounders).} \\
\\ 
\texttt{1.  **The Causal Structure**:} \\
\texttt{    *   \$U \textbackslash{}rightarrow \textbackslash{}text\{Education\}\$} \\
\texttt{    *   \$U \textbackslash{}rightarrow \textbackslash{}text\{Salary\}\$} \\
\texttt{    *   \$\textbackslash{}text\{Proximity\} \textbackslash{}rightarrow \textbackslash{}text\{Education\}\$} \\
\texttt{    *   \$\textbackslash{}text\{Education\} \textbackslash{}rightarrow \textbackslash{}text\{Salary\}\$} \\
\texttt{    *   \$U\$ is unobserved, meaning we cannot condition on it to block the backdoor path between Education and Salary.} \\
\\ 
\texttt{2.  **Analyzing the Data via Proximity as a Proxy**:} \\
\texttt{    Since \$U\$ is unobserved, we cannot calculate the causal effect of Education on Salary directly using standard adjustment methods on the whole population. However, the problem provides data stratified by **Proximity to a college**.} \\
\texttt{    *   **Crucial Causal Assumption**: The prompt states "Proximity to a college has a direct effect on education level." It does *not* state that Proximity has a direct effect on Salary, nor does it state that Proximity is a descendant of \$U\$. In standard causal diagrams for this type of problem} \\
            }
        }
    }
    & & \\ \\

    \theutterance \stepcounter{utterance}  
    & & & \multicolumn{2}{p{0.3\linewidth}}{
        \cellcolor[rgb]{0.9,0.9,0.9}{
            \makecell[{{p{\linewidth}}}]{
                \texttt{\tiny{[GM$|$GM]}}
                \texttt{No} \\
            }
        }
    }
    & & \\ \\

    \theutterance \stepcounter{utterance}  
    & & & \multicolumn{2}{p{0.3\linewidth}}{
        \cellcolor[rgb]{0.9,0.9,0.9}{
            \makecell[{{p{\linewidth}}}]{
                \texttt{\tiny{[GM$|$GM]}}
                \texttt{no} \\
            }
        }
    }
    & & \\ \\

    \theutterance \stepcounter{utterance}  
    & & & \multicolumn{2}{p{0.3\linewidth}}{
        \cellcolor[rgb]{0.9,0.9,0.9}{
            \makecell[{{p{\linewidth}}}]{
                \texttt{\tiny{[GM$|$GM]}}
                \texttt{game\_result = WIN} \\
            }
        }
    }
    & & \\ \\

\end{supertabular}
}

\end{document}
